% Chapter Template

\chapter{Conclusiones} % Main chapter title

\label{Chapter5} % Change X to a consecutive number; for referencing this chapter elsewhere, use \ref{ChapterX}
En este capítulo se presentan las conclusiones obtenidas de la ejecución de este
trabajo y, además, se proponen perspectivas que podrían ser implementadas a
futuro.


%----------------------------------------------------------------------------------------

%----------------------------------------------------------------------------------------
%	SECTION 1
%----------------------------------------------------------------------------------------

\section{Conclusiones generales }


En primer lugar, se destaca que fue posible desarrollar con éxito un
dispositivo capaz de comunicarse tanto con la red de domótica de
Enertik como con el ecosistema Victron.

Este resultado resulta de gran relevancia, ya que habilita la futura
integración de equipos Victron dentro de la red de domótica del
motorhome, amplía las capacidades del sistema y permite una mayor
flexibilidad en su implementación.

Asimismo, se logró el desarrollo de una aplicación funcional que
permite interactuar con el sistema de domótica. Esta aplicación es
capaz de replicar la mayoría de las funcionalidades disponibles en la
pantalla HMI del motorhome y brinda a los usuarios una alternativa
adicional para el control del sistema.

A continuación, se destacan los principales aspectos del trabajo:

\begin{itemize}
\item Se desarrolló con éxito un dispositivo capaz de comunicarse con ambos ecosistemas de manera confiable.
\item Se diseñó un PCB de tamaño reducido, orientado a su futura implementación como producto comercial.
\item El equipo se integra a la red del motorhome mediante el cableado estándar de cuatro conductores, lo que facilita su instalación.
\item Se implementó una arquitectura de firmware basada en múltiples tareas, lo que permite un funcionamiento organizado y eficiente del sistema.
\item La aplicación desarrollada funciona correctamente en dispositivos Android, lo que amplía las posibilidades de uso por parte de los usuarios.
\item La aplicación replica la mayoría de las funcionalidades de la pantalla HMI, lo que reduce la dependencia de esta para el control del sistema.
\item Se validó el funcionamiento del sistema tanto en entorno de laboratorio como en condiciones reales de operación.
\item La solución presenta un diseño escalable, lo que permite la incorporación de nuevas funcionalidades y dispositivos en el futuro.
\item Se emplearon protocolos de comunicación estándar, lo que favorece la interoperabilidad con otros sistemas.
\end{itemize}


%----------------------------------------------------------------------------------------
%	SECTION 2
%----------------------------------------------------------------------------------------
\section{Próximos pasos}

Como líneas de trabajo a futuro, se consideran diversas mejoras y
extensiones del sistema desarrollado.

En primer lugar, durante los ensayos en vehículos reales, los
fabricantes de motorhomes señalaron que la aplicación requiere la
conexión a la red Wi-Fi generada por el ESP32, lo que implica que el
teléfono celular pierde el acceso a Internet mientras se utiliza el
sistema. Si bien esta situación no resulta crítica, se analizaron
distintas alternativas para su resolución.

Una posible solución consiste en modificar el firmware del ESP32 para
que se conecte a una red Wi-Fi existente. De esta manera, la aplicación,
al conectarse a la misma red, podría acceder a los \textit{endpoints}
del servidor mediante la dirección IP asignada por el módem.

Otra alternativa consiste en migrar el protocolo de comunicación entre
la aplicación y el ESP32. En este sentido, se podría implementar una
comunicación mediante Bluetooth en reemplazo de la conexión Wi-Fi. Ambas
opciones permitirían mantener la conectividad a Internet en el
dispositivo móvil y, al mismo tiempo, conservar el control del sistema
sin inconvenientes.

Por otra parte, se considera necesaria la adaptación de la aplicación
para su ejecución en dispositivos iOS. Esta tarea no se incluyó dentro
del alcance del presente trabajo. 

En la actualidad, los sistemas
operativos predominantes en dispositivos móviles son Android e iOS. La compatibilidad con ambos sistemas permitiría ampliar el alcance de la solución, al cubrir una mayor cantidad de usuarios potenciales.

Relacionado con este aspecto, también resulta necesario realizar
modificaciones que permitan distribuir la aplicación a través de
tiendas oficiales, como Play Store en Android y App Store en iOS.

En la actualidad, la aplicación solo puede instalarse mediante un
archivo de tipo APK. Sin embargo, para su implementación a nivel
comercial, resulta conveniente que los usuarios puedan acceder a ella a
través de estas plataformas, lo que facilitaría su distribución,
instalación y actualización.